\section{Objective identities and conditional bounds}
\label{app:theory}

Let $H\in\{0,1\}^{n\times k}$ be the assignment matrix, with one nonzero per row,
and let $h_c$ denote its columns. The cross-fold mask is
$\mathbf{1}\mathbf{1}^\top-HH^\top$ and has zero diagonal. Consequently,
\begin{equation}
 \cut_S(H)=\tfrac12\left(\mathbf{1}^\top S\mathbf{1}
                       -\sum_c h_c^\top S h_c\right).
\end{equation}
Substituting $S=BB^\top$ yields Equation~\ref{main-eq:factor}. Group contraction is exact
because a whole group's contribution to every block sum is the sum of its rows.
This statement concerns retained groups; dissolving a group changes the admissible
partition set rather than invalidating the identity.

For a positive-semidefinite kernel with features $\phi_i$, its integer elementwise
power has features $\phi_i^{\otimes p}$ because
$\langle\phi_i^{\otimes p},\phi_j^{\otimes p}\rangle=
\langle\phi_i,\phi_j\rangle^p$. This gives the stated positive-semidefinite
property directly and explains why arbitrary fractional powers are not accepted.

For nonnegative similarities, let $N_\tau(z)$ count unordered crossing pairs
with $S_{ij}\ge\tau>0$. Each such pair contributes at least $\tau^p$, so
\begin{equation}
 \cut_{S^{\circ p}}(z)\ge\tau^p N_\tau(z).
\end{equation}
Thus the powered cut supplies a soft penalty on the number of strong crossing
pairs. For a fixed nonempty set of crossing pairs,
$\lim_{p\to\infty}\cut_{S^{\circ p}}(z)^{1/p}
=\max_{i<j:z_i\ne z_j}S_{ij}$, by the finite-dimensional $\ell_p$ norm limit.
This is the maximum over all crossing pairs, not the average of per-test-point
nearest neighbours used in $\nnl$. The power therefore provides a principled
tail surrogate, without making the optimized objective identical to the
evaluation metric. A rank-limited factor introduces a separate approximation
error for each power.

For symmetric residual $E=S-BB^\top$, the spectral norm inequality
$|v^\top E v|\le\|E\|_2\|v\|^2$ gives
\begin{align}
 |\cut_S(H)-\cut_{BB^\top}(H)|
 &\le\tfrac12\|E\|_2\left(\|\mathbf{1}\|^2+\sum_c\|h_c\|^2\right)\\
 &=n\|E\|_2\equiv\delta.
 \label{eq:bound}
\end{align}
The bound is uniform across assignments and can be loose in practice. If a solver
returns $\widehat H$ with approximate-kernel objective at most $\eta$ above the
minimum over a specified feasible set, then
\begin{equation}
 \cut_S(\widehat H)\le\min_{H\ \mathrm{feasible}}\cut_S(H)+2\delta+\eta.
\end{equation}
Our heuristic does not certify $\eta$, so this conditional inequality is not an
empirical optimality guarantee. It also requires comparing the same feasible set;
grouping, count changes and kernel sharpening alter the optimization problem.

Moving a factor row $b$ from block $a$ to $c$ replaces $(p_a,p_c)$ by
$(p_a-b,p_c+b)$. Expansion of Equation~\ref{main-eq:factor} gives
$\Delta\cut=b^\top p_a-b^\top p_c-\|b\|^2$. For a swap of factor rows $u$ and $v$,
the reduction equals the sum of their single-move reductions minus $2u^\top v$.
These identities hold for group sums as well. Searching only a finite set of swap
candidates preserves the validity of an accepted move, but does not establish
that no improving swap exists elsewhere.

If all fold counts equal a fixed value $m$, kernel k-means maximizes
$m^{-1}\sum_c\|p_c\|^2$ and agrees with dense-cut minimization. For unequal fixed
counts it maximizes a differently weighted sum, and for a capacity corridor the
weights themselves change with the assignment. This is why we do not transfer
ordinary kernel k-means convergence or approximation guarantees to the present
solver without further assumptions.

